\section*{The Appointed Propers in Full}

\noindent The Latin below is collated against the 1962 Vatican typical Missal,
printed pp.~389--390. Scriptural English is the public-domain
Douay--Rheims/Challoner; the orations are quoted from the public-domain 1861
Cummiskey hand missal. The English witnesses aid study and are not represented
as the official 1962 liturgical text.

\fulltextelement{Introit\enspace\properrefs{Int.}}{\latin{Psalmus 54:17, 18, 20, 23; 54:2}}
Cum clamárem ad Dóminum, exaudívit vocem meam, ab his qui appropínquant mihi:
et humiliávit eos qui est ante s\'{\ae}cula, et manet in ætérnum: iacta cogitátum
tuum in Dómino, et ipse te enútriet.

\medskip\noindent\latin{Psalmus.} Exáudi, Deus, oratiónem meam, et ne despéxeris
deprecatiónem meam: inténde mihi, et exáudi me.
\finishfulltextelement
\begin{fulltextenglish}{Douay--Rheims/Challoner, Ps.~54:17--18, 20, 23, 2}
But I have cried to God: and the Lord will save me. Evening and morning, and
at noon I will speak and declare: and he shall hear my voice. God shall hear,
and the Eternal shall humble them. Cast thy care upon the Lord, and he shall
sustain thee. Hear, O God, my prayer, and despise not my supplication.
\end{fulltextenglish}

\fulltextelement{Collect\enspace\properrefs{Coll.}}{\latin{Oratio}}
Deus, qui omnipoténtiam tuam parcéndo máxime et miserándo maniféstas:
multíplica super nos misericórdiam tuam; ut, ad tua promíssa curréntes,
cæléstium bonórum fácias esse consórtes. Per Dóminum nostrum Iesum Christum,
Filium tuum: Qui tecum vivit et regnat in unitate Spiritus Sancti Deus, per
omnia saecula saeculorum.
\finishfulltextelement
\begin{fulltextenglish}{Cummiskey hand missal (1861), Collect, p.~416}
O God, who chiefly manifestest thy Almighty power in pardoning and shewing
mercy, increase thy goodness towards us: that having recourse to thy promises,
we may be partakers of thy heavenly blessing. Thro'.
\end{fulltextenglish}

\fulltextelement{Epistle\enspace\properrefs{Ep.}}{\latin{1 ad Corinthios 12:2--11}}
Fratres: Scitis quóniam cum gentes essétis, ad simulácra muta prout ducebámini
eúntes. Ideo notum vobis fácio, quod nemo in Spíritu Dei loquens, dicit
anáthema Iesu. Et nemo potest dícere, Dóminus Iesus, nisi in Spíritu Sancto.
Divisiónes vero gratiárum sunt, idem autem Spíritus. Et divisiónes
ministratiónum sunt, idem autem Dóminus. Et divisiónes operatiónum sunt, idem
vero Deus, qui operátur ómnia in ómnibus. Unicuíque autem datur manifestátio
Spíritus ad utilitátem. Alii quidem per Spíritum datur sermo sapiéntiæ: álii
autem sermo sciéntiæ secúndum eúndem Spíritum: álteri fides in eódem Spíritu:
álii grátia sanitátum in uno Spíritu: álii operátio virtútum, álii prophétia,
álii discrétio spirítuum, álii génera linguárum, álii interpretátio sermónum.
Hæc autem ómnia operátur unus atque idem Spíritus, dívidens síngulis prout
vult.
\finishfulltextelement
\begin{fulltextenglish}{Douay--Rheims/Challoner, 1~Cor.~12:2--11}
You know that when you were heathens, you went to dumb idols, according as you
were led. Wherefore, I give you to understand that no man, speaking by the
Spirit of God, saith Anathema to Jesus. And no man can say The Lord Jesus, but
by the Holy Ghost. Now there are diversities of graces, but the same Spirit.
And there are diversities of ministries, but the same Lord. And there are
diversities of operations, but the same God, who worketh all in all. And the
manifestation of the Spirit is given to every man unto profit. To one indeed,
by the Spirit, is given the word of wisdom: and to another, the word of
knowledge, according to the same Spirit: to another, faith in the same spirit:
to another, the grace of healing in one Spirit: to another the working of
miracles: to another, prophecy: to another, the discerning of spirits: to
another, diverse kinds of tongues: to another, interpretation of speeches. But
all these things, one and the same Spirit worketh, dividing to every one
according as he will.
\end{fulltextenglish}

\fulltextelement{Gradual and Alleluia\enspace\properrefs{Grad., All.}}{\latin{Psalmi 16:8, 2; 64:2}}
Custódi me, Dómine, ut pupíllam óculi: sub umbra alárum tuárum prótege me.
\medskip\noindent\latin{℣.} De vultu tuo iudícium meum pródeat: óculi tui
vídeant æquitátem.

\medskip\noindent Allelúia, allelúia. \latin{℣.} Te decet hymnus, Deus, in
Sion: et tibi reddétur votum in Ierúsalem. Allelúia.
\finishfulltextelement
\begin{fulltextenglish}{Douay--Rheims/Challoner, Ps.~16:8, 2; 64:2}
From them that resist thy right hand keep me, as the apple of thy eye. Protect
me under the shadow of thy wings. Let my judgment come forth from thy
countenance: let thy eyes behold the things that are equitable. A hymn, O God,
becometh thee in Sion: and a vow shall be paid to thee in Jerusalem.
\end{fulltextenglish}

\fulltextelement{Gospel\enspace\properrefs{Gosp.}}{\latin{Lucas 18:9--14}}
In illo témpore: Dixit Iesus ad quosdam, qui in se confidébant tamquam iusti,
et aspernabántur céteros, parábolam istam: Duo hómines ascendérunt in templum
ut orárent: unus pharis\'{\ae}us, et alter publicánus. Pharis\'{\ae}us stans, hæc apud
se orábat: Deus, grátias ago tibi, quia non sum sicut céteri hóminum:
raptóres, iniústi, adúlteri: velut étiam hic publicánus. Ieiúno bis in sábbato:
décimas do ómnium, quæ possídeo. Et publicánus a longe stans nolébat nec
óculos ad cælum leváre: sed percutiébat pectus suum, dicens: Deus, propítius
esto mihi peccatóri. Dico vobis: descéndit hic iustificátus in domum suam ab
illo: quia omnis qui se exáltat, humiliábitur: et qui se humíliat, exaltábitur.
\finishfulltextelement
\begin{fulltextenglish}{Douay--Rheims/Challoner, Luke 18:9--14}
And to some who trusted in themselves as just and despised others, he spoke
also this parable: Two men went up into the temple to pray: the one a Pharisee
and the other a publican. The Pharisee standing, prayed thus with himself: O
God, I give thee thanks that I am not as the rest of men, extortioners, unjust,
adulterers, as also is this publican. I fast twice in a week: I give tithes of
all that I possess. And the publican, standing afar off, would not so much as
lift up his eyes towards heaven; but struck his breast, saying: O God, be
merciful to me a sinner. I say to you, this man went down into his house
justified rather than the other: because every one that exalteth himself shall
be humbled: and he that humbleth himself shall be exalted.
\end{fulltextenglish}

\fulltextelement{Offertory\enspace\properrefs{Off.}}{\latin{Psalmus 24:1--3}}
Ad te, Dómine, levávi ánimam meam: Deus meus, in te confído, non erubéscam:
neque irrídeant me inimíci mei: étenim univérsi, qui te exspéctant, non
confundéntur.
\finishfulltextelement
\begin{fulltextenglish}{Douay--Rheims/Challoner, Ps.~24:1--3}
To thee, O Lord, have I lifted up my soul. In thee, O my God, I put my trust;
let me not be ashamed. Neither let my enemies laugh at me: for none of them
that wait on thee shall be confounded.
\end{fulltextenglish}

\fulltextelement{Secret\enspace\properrefs{Sec.}}{\latin{Secreta; Praefatio de Sanctissima Trinitate}}
Tibi, Dómine, sacrifícia dicáta reddántur: quæ sic ad honórem nóminis tui
deferénda tribuísti, ut eádem remédia fíeri nostra præstáres. Per Dóminum
nostrum Iesum Christum, Filium tuum: Qui tecum vivit et regnat in unitate.
\finishfulltextelement
\begin{fulltextenglish}{Cummiskey hand missal (1861), Secret, p.~417}
May the sacrifice, we offer, O Lord, be presented before thee: which thou hast
appointed to be offered in honour of thy name; and at the same time become a
remedy to us. Thro'.
\end{fulltextenglish}

\fulltextelement{Communion\enspace\properrefs{Comm.}}{\latin{Psalmus 50:21}}
Acceptábis sacrifícium iustítiæ, oblatiónes et holocáusta, super altáre tuum,
Dómine.
\finishfulltextelement
\begin{fulltextenglish}{Douay--Rheims/Challoner, Ps.~50:21}
Then shalt thou accept the sacrifice of justice, oblations and whole burnt
offerings: then shall they lay calves upon thy altar.
\end{fulltextenglish}

\fulltextelement{Postcommunion\enspace\properrefs{Postcomm.}}{\latin{Postcommunio}}
Qu\'{\ae}sumus, Dómine Deus noster: ut, quos divínis reparáre non désinis
sacraméntis, tuis non destítuas benígnus auxíliis. Per Dóminum nostrum.
\finishfulltextelement
\begin{fulltextenglish}{Cummiskey hand missal (1861), Postcommunion, p.~418}
We beseech thee, O Lord our God, that in thy mercy thou wouldst never deprive
those of thy help, whom thou continually strengthenest by these divine
mysteries. Thro'.
\end{fulltextenglish}
